\section{Calendar order}
\label{sec:calendar}

Butterfly freedom is a property of one marginal law.  A surface also requires
later expiries to dominate earlier ones in convex order.  Under deterministic
rates and carry, this comparison has a simple form in normalized units.

\begin{lemma}[Calendar ordering at equal log-moneyness]
\label{lem:caljensen}
Let $F_{0,t}$ be the deterministic forward curve and suppose
$M_t=S_t/F_{0,t}$ is a positive martingale.  For $T_1<T_2$ and every $k$,
the normalized calls satisfy
\begin{equation}
c_1(k)\le c_2(k).
\label{eq:calreq}
\end{equation}
\end{lemma}

\begin{proof}
Since $M_{T_i}=S_{T_i}/F_i=Y_i$, conditional Jensen gives
\[
\begin{aligned}
c_2(k)
&=\E[(M_{T_2}-e^k)^+]\\
&\ge\E[(\E[M_{T_2}\mid\mathcal F_{T_1}]-e^k)^+]
=\E[(M_{T_1}-e^k)^+]=c_1(k).
\end{aligned}
\]
\end{proof}

Equal $k$ means strikes $K_i=F_i e^k$, which are generally different dollar
strikes.  The comparison is a forward-scaled calendar relation, not the usual
listed equal-strike spread.  In probability language, the laws of $Y_i$ have
equal means and increase in convex order.

The scaling is essential.  A violation at $k$ is monetized by shorting
$1/(D_1F_1)$ calls at $K_1=F_1e^k$ and buying $1/(D_2F_2)$ calls at
$K_2=F_2e^k$.  After normalization both payoffs are convex functions of the
common martingale $M_t$.  Without deterministic carry, one must first specify
the numeraire and the objects being compared; equal log-moneyness alone is not
a model-free calendar relation.

\subsection{Calls and ledgers are conjugate}

\begin{lemma}[Ledger--call conjugacy]
\label{lem:conjugacy}
For any admissible slice,
\begin{equation}
\begin{aligned}
c(k)&=\max_{z\in[-\infty,\infty]}
\{G(z)-e^k(1-\Lambda(z))\},\\
G(z)&=\min_{k\in[-\infty,\infty]}
\{c(k)+e^k(1-\Lambda(z))\}.
\end{aligned}
\label{eq:conjugacy}
\end{equation}
The optimizers are $z=z_k$ and $k=x(z)$, respectively.
\end{lemma}

\begin{proof}
For fixed $k$,
\[
G(z)-e^k(1-\Lambda(z))
=\int_z^\infty(e^{x(t)}-e^k)\rho(t)\,\dd t.
\]
The integrand is negative before $z_k$ and positive after it.  Starting the
integral at $z_k$ therefore maximizes it and gives the call price.  The first
identity implies
$c(k)+e^k(1-\Lambda(z))\ge G(z)$ for every $k$, with equality at
$k=x(z)$, proving the second.
\end{proof}

Geometrically, the first identity scans all possible exercise ranks.  Before
$z_k$, the integral includes outcomes where the payoff is negative; after
$z_k$, it discards outcomes where the payoff is positive.  The strike root is
the unique point of tangency between call and ledger descriptions.  The
second identity reconstructs the ledger from the full call curve, so no
information is lost in changing representation.

\begin{theorem}[Calendar order is ledger order]
\label{thm:ledgerorder}
For two admissible mean-one slices,
\begin{equation}
c_1(k)\le c_2(k)\ \text{for all }k
\quad\Longleftrightarrow\quad
G_1(z)\le G_2(z)\ \text{for all }z.
\label{eq:ledgerorder}
\end{equation}
\end{theorem}

\begin{proof}
If $G_1\le G_2$, the maximum representation in \eqref{eq:conjugacy} gives
$c_1\le c_2$.  If $c_1\le c_2$, the minimum representation gives
$G_1\le G_2$.  Both implications are pointwise.
\end{proof}

Thus the continuum of calendar inequalities is equivalent to ordering two
arrays already produced by the slice builds.  This is also the classical
convex-order criterion: by the results of Strassen and Kellerer, a family of
equal-mean laws increasing in convex order can be realized as the marginals of
a martingale \citep{Strassen1965,Kellerer1972}.

This is an existence statement.  Ordered ledgers make the marginal laws
compatible with some martingale dynamics; they neither select that dynamics
nor determine forward smiles.  The calendar problem treated here is exactly
the static marginal problem.

\subsection{A support-confined control}
\label{sec:calcontrol}

The theorem is an exact diagnostic.  A finite calibration requires a chosen
enforcement rule.  When fitting a farther expiry after a nearer one, the
reference objective adds
\begin{equation}
\sqrt{w_{\rm cal}}\,\zeta_m
\{c_{\rm near}(k_m)-c_{\rm far}(k_m)\}^+
\label{eq:calrows}
\end{equation}
on a grid of strikes in the intersection of their quoted spans.  The weights
$\zeta_m$ taper to zero with a $\cos^2$ profile over the outer fifteen percent
of that span.  The number of nodes is $\max(49,4N+1)$.

The rows are imposed in price space rather than ledger space.  A ledger value
at one rank integrates the entire upper tail, so constraining it makes one
expiry's unquoted tail move another.  In a reproduced stress comparison, a
full-grid ledger floor moved the worst quote error from
\MacPhantomDragFromBp\ to \MacPhantomDragToBp\ vol bp.  Fixed-strike prices
localize the constraint to the region being compared.

The weight is finite.  Any remaining violation is therefore reported as a
price gap and as a fraction of the local bid--ask width.  The exact scope is:
calendar order is controlled on common quoted support; outside that support,
the full ledger comparison remains a diagnostic of the extrapolations.

The theorem itself is not imposed at every rank because $G(z)$ is a tail
integral.  Even at a central rank it contains every more extreme outcome, so
a constraint there couples quoted prices to the entire unquoted continuation.
The ledger is ideal for diagnosis precisely because it sees the whole law;
fixed-strike call rows are better for calibration because each corresponds to
a specified payoff in the common quoted region.

\begin{figure}[htbp]
\centering
\paperfig{0.92\textwidth}{fig_calendar.pdf}
\caption{Calendar order in the frozen SPY sample.  Left: total variance is
ordered through maturity at representative log-moneyness values.  Right: the
two shortest expiries are ordered on their common quoted span, while a
volatility-space crossing appears only in a numerically empty far wing.  The
price-space gap there is at round-off scale.}
\label{fig:calendar}
\end{figure}

For the two shortest SPY expiries, a full-grid volatility audit flags a
\MacCalGapVolBp\ vol-bp crossing at $k=\MacCalGapArgmaxK$, far outside the
common quote span $[\MacCalSpanLo,\MacCalSpanHi]$.  The corresponding calls
are \MacCalendarCallShort\ and \MacCalendarCallLong\ of forward.  A direct
price audit finds a worst violation of only \MacCalPriceGridWorst.  The
crossing is a disagreement between extrapolations in a region where Black
inversion is ill-conditioned, not an economically available calendar trade.
